% META: {"id":"1SPE-GEOREP-EX-048","chapitre":"1SPE-GEOMETRIE-REPEREE","type_objet":"exercice","capacites":["1SPE-GEOMETRIE-REPEREE-C5"],"capacites_codes":["C5"],"methodes":["M5"],"parcours":3,"competences":["calculer","raisonner"],"duree_min":20,"mode_creation":"ex_nihilo","sources_inspiration":[],"fichier_tex":"chapitres/1SPE-GEOMETRIE-REPEREE/exercices/1SPE-GEOREP-EX-048.tex","corrige_tex":"chapitres/1SPE-GEOMETRIE-REPEREE/corriges/1SPE-GEOREP-CO-048.tex","status":"approved","gestes":["raisonnement_par_coordonnees"],"origin":"generated","status_history":["generated","audited","accepted","approved"],"release_acceptance":"RELEASE_OWNER_BATCH_ACCEPTANCE","acceptance_closure_digest":"sha256:9b3ccf9a81c5520fbb7e03b7b2d3e7bf2057a02834b6d908bf3be5f49f3b553f","acceptance_receipt_digest":"sha256:4fad86f0282e0fa4e0419cca1ad6379ae7af5a280c04a4a90880f90a1c044242"}
% BEGIN-VERIFY
% from sympy import *
% x, y = symbols('x y')
% # Cercle circonscrit au triangle A(0,0), B(4,0), C(0,6)
% # x^2+y^2+Ax+By+C=0. A: C=0. B: 16+4A=0 => A=-4. C: 36+6B=0 => B=-6
% # Centre (2,3), rayon sqrt(4+9)=sqrt(13)
% assert (4+9) == 13
% END-VERIFY
\begin{exercice}{1SPE-GEOREP-EX-048}{3}{20}
Déterminer l'équation du cercle circonscrit au triangle $ABC$ avec $A(0 ; 0)$, $B(4 ; 0)$ et $C(0 ; 6)$.
\end{exercice}
